\subsection{An aligned bit circuit with cheaper centers}
\label{sec:aligned-bit}
We change the bit network of Section~\ref{sec:paired-bit-construction} in
two ways: a smaller side circuit, and center wires that lose $h-1$ rather
than $h$ dimensions per invocation. The ground size $h=50$, the form
$I-J/9$, the stage frames of the data banks, and the bank sharing between
stages 1 and 3 are unchanged.

\paragraph{Aligned blocks.} In the weighted paired recursion for the common
point $i$, put the points $2k,2k+1$ in one block whenever both differ from
$i$; the partner $i'$ of $i$ forms a singleton. At every coarser level group
blocks $2k,2k+1$ in the same way. All groups therefore use the same blocks
away from their own point, and the identification of equal pair-star sums
across groups applies more often.

\paragraph{Shared pair-star sums.} At the top level of group $i$, a vertex
$a$ whose partner $q$ lies in its block $A$ needs, for each other block $J$,
the sum $Y_i(a,J)$ of the triples $\{i,q,c\}$ with $c\notin A\cup J$. Let
$\mathcal C$ consist of the blocks other than $A$ and the block of $i$. Form
the prefix and suffix sums, in increasing block order over $\mathcal C$, of
the pair-star block sums $\sum_{c\in C}x_{iqc}$. These chains depend only on
the unordered pair $\{i,q\}$, since $\mathcal C$ consists of the full blocks
avoiding both $i$ and $q$; group $q$ forms the same chains for its vertex
$i'$. For $J\in\mathcal C$, $Y_i(a,J)$ is the prefix--suffix sum plus the
single triple $\{i,q,i'\}$; for $J$ the block of $i$ it is the total over
$\mathcal C$. The rest of the recursion is unchanged.

\paragraph{Totals.} Each group also keeps its total
$z_i=\sum_{T\ni i}x_T$, the top-level total of its recursion. Its source span
is $H_i=\langle t_T:T\ni i\rangle$, which is positive definite by the
common-point argument above.

Equal supports are identified across groups as before, by the fixed pair and
the third points of a pair-star. At $h=50$ the retained graph has
$c=435346$ additions, the $q=58800$ partial outputs, and $h=50$ totals.
Exact support comparison checks every partial output and total, the
disjointness of every addition, and the common point of every node.

\paragraph{Centers from the totals.} Let the center wire $C_i$ be the role
carrying $z_i$ out of the side graph. Compile roles as before, treating the
$h$ center uses like designated outputs, so $R=c+q+h=494196$ roles suffice.
Let $L'$ be the product of the gates of this graph, $V$ the source copy, $J$
the injection of the partial outputs, $R_0$ the original central scatter,
and $J'=JR_0$. The forward schedule is
\[
 L',\ J',\ L'^{-1},\ V,\ L',\ J',\ L'^{-1},\ V^{-1}.
\]
For arbitrary scratch $z$ the injections give $J'L'z+J'L'(z+Vx)=J'L'Vx$.
The center roles of $L'Vx$ carry $z_i=(Gx)_i$, so $J'L'V=JLV+R_0G=I$ and the
invocation is $y\gets y+x$; the cleanup restores every auxiliary value.
Stage 2 reverses and inverts this schedule.

In a forward invocation the frames are $D_0$ for the first three operations,
$X_{t_X}$ for $V$, the graded frames $D_{U_z}$ for the middle $L'$, $D_0$ for
$R_0$, $Y_{t_Y}$ for $J$, and $D_1$ for the last two operations. Thus $C_i$
passes through $D_{H_i}$ at its gate in the middle $L'$ and returns to $D_0$
at $R_0$, losing $P\otimes H_i$, of dimension $h-1$. The physical $Y$ wires
pass $B\otimes\langle t\rangle\subset D_0\subset Y_t$, and the physical $X$
wires pass $A\otimes\langle t\rangle\subset X_t\subset D_1$.

In stage 2 the operations are
\[
 V,\ L',\ J,\ R_0,\ L'^{-1},\ V^{-1},\ L',\ J,\ R_0,\ L'^{-1},
\]
with logical source $Y$ and target $X$. Their frames are $D_0$, $D_0$,
$X_{t_X}$, $D_1$, the complement frames $D_{U_z^\perp}$, $Y_{t_Y}$, and then
$D_1$. The center $C_i$ rises from $D_0$ to $D_1$ at the first $R_0$ and
falls to $D_{H_i^\perp}$ at its gate in $L'^{-1}$, again losing
$P\otimes H_i$. Since $H_i$ is positive definite, $H_i^\perp$ is
nondegenerate. The physical $X$ wires pass
$A\otimes\langle t\rangle\subset X_t\subset D_1$, and the partial-output
roles pass $\langle t_X\rangle\subset U_z^\perp$ as before. Every other edge is
one of the nested edges already treated. Each auxiliary role starts at $D_0$
and ends at $D_1$, so the stage-1/stage-3 bank sharing applies unchanged.

\begin{proposition}[Aligned bit interface]\label{prop:aligned-bit-interface}
The aligned bit network restores arbitrary scratch, exchanges the data banks,
and satisfies the original all-role endpoint identity with
\begin{align*}
 W_{\rm b}^{\rm al}&=2N+2v^2R=394759742720000,\\
 s_{\rm b}^{\rm al}&=W_{\rm b}^{\rm al}m-N+6v^2h(h-1)
                         =49344965957616000000,\\
 W_{\rm b}^{\rm al}m-s_{\rm b}^{\rm al}&=1882384000000>0.
\end{align*}
\end{proposition}

\begin{proof}
The scalar schedule and both frame assignments were given above. The only
decreasing edges are the $3v^2h$ center returns, each of dimension $h-1$.
Telescoping gives total absolute label change
$W_{\rm b}^{\rm al}m-2N+6v^2h(h-1)$, and the unchanged negative source
projections add $N$ rational rank. Data endpoints are unchanged and every
auxiliary role has endpoints zero and $I_m$, with its scalar value restored.
The common-frame identity and rational-shear transfer apply to these fixed
finite data.
\end{proof}

The exact deficit is $\eta_{\rm b}^{\rm al}=49/1284490000$, against
$23/661055000$ for the paired network. With $\log m<11737/1000$,
$\eta_{\rm b}^{\rm al}>(325/10^{11})(11737/1000)$, so
$s_{\rm b}^{\rm al}/W_{\rm b}^{\rm al}<m^{1-325/10^{11}}$.
